\subsection{Случай одной независимой переменной}

\begin{definition}
Функция $u : (a,b) \to \R$ называется \textbf{абсолютно непрерывной}, если
\[
  \exists\, v \in \Lp{1}(a,b): \qquad u(x) = \int\limits_a^x v(t)\, dt + C .
\]
\end{definition}

\begin{remark}
$u$ абсолютно непрерывна $\Longrightarrow$ $u$ непрерывна.
\end{remark}

\begin{warn}
В тетради теорема ниже сначала была записана в виде «для того, чтобы $u \in \Lp{1}(a,b)$
имела первую обобщённую производную, необходимо и достаточно, чтобы $u$ была абсолютно
непрерывна», и её доказательство занимало две страницы — \textbf{обе перечёркнуты
крест-накрест}. Затем формулировка и доказательство переписаны заново; ниже приведён
именно второй, итоговый вариант. В зачёркнутой версии доказательство
необходимости шло через $W(x) = \int_a^x \frac{du}{dt}(t)\,dt$ и равенство
$\frac{d(W-u)}{dx} = 0 \Rightarrow W = u + C$ п.\,в.
\end{warn}

\begin{theorem}\label{thm:ac}
Функция $u \in \Lp{1}(a,b)$ абсолютно непрерывна $\iff$ существует
$u' \in \Lp{1}(a,b)$ (здесь $u'$ — соболевская производная).
\end{theorem}

\begin{proof}
\textbf{$\Longrightarrow$.} Пусть $u \in \Lp{1}(a,b)$ абсолютно непрерывна, то есть
\[
  \exists\, v \in \Lp{1}(a,b): \quad u(x) = \int\limits_a^x v(t)\, dt + C .
\]
Пользуемся признаком существования производной (теорема~\ref{thm:crit-deriv}).
Гладкие функции плотны в $\Lp{1}$, поэтому
\[
  \exists\, v_m \in C^\infty\big((a,b)\big): \quad v_m \to v \ \text{ в } \Lp{1}(a,b) .
\]
Положим
\[
  u_m(x) = \int\limits_a^x v_m(t)\, dt + C .
\]
Тогда для всех $x \in [a,b]$
\[
  |u_m(x) - u(x)| \le \int\limits_a^x |v_m(t) - v(t)|\, dt
  \le \int\limits_a^b |v_m(t) - v(t)|\, dt = \norm{v_m - v}_{\Lp{1}(a,b)} ,
\]
то есть $u_m \unif u$ на $(a,b)$, а значит $u_m \to u$ в $\Lp{1}(a,b)$.

Так как $u_m, v_m \in C^\infty$, то $u_m' = v_m$, причём
\[
  u_m \to u, \qquad v_m \to v \qquad \text{в } \Lp{1}(a,b) .
\]
Тогда по признаку существования производной $u' = v \in \Lp{1}(a,b)$.

\textbf{$\Longleftarrow$.} Пусть $u \in \Lp{1}(a,b)$ и существует $u' = w \in \Lp{1}(a,b)$.
Положим
\[
  \widetilde u(x) = \int\limits_a^x w(t)\, dt + C ,
\]
тогда $\widetilde u$ абсолютно непрерывна, и по уже доказанной части
$\widetilde u\,' = w$. Отсюда
\[
  (u - \widetilde u)' = w - w = 0
  \quad\overset{\text{Л\ref{lem:const}}}{\Longrightarrow}\quad
  u = \widetilde u + C_1
  \quad\Longrightarrow\quad
  u(x) = \int\limits_a^x w(t)\, dt + \underbrace{C + C_1}_{C_2} ,
\]
то есть $u$ абсолютно непрерывна.
\end{proof}

\subsubsection*{Почему это называется абсолютной непрерывностью}

Если $u \in C\big([a,b]\big)$, то $u$ равномерно непрерывна:
\[
  \forall \eps > 0 \ \exists \delta > 0: \quad |x - y| < \delta
  \ \Longrightarrow\ |u(x) - u(y)| < \eps .
\]
Классическое определение абсолютной непрерывности требует большего — контроля не над
одним отрезком, а над любым конечным набором непересекающихся отрезков:
\begin{idea}[title={Классическое определение}]
Для $a \le x_1 < y_1 < x_2 < y_2 < \dots < x_k \le y_k \le b$
\[
  \sum_{i=1}^{k} (x_i - y_i) < \delta
  \quad\Longrightarrow\quad
  \sum_{i=1}^{k} |u(x_i) - u(y_i)| < \eps .
\]
\end{idea}

\begin{warn}
В тетради в сумме написано $\sum (x_i - y_i)$ при $x_i < y_i$, то есть формально
отрицательные слагаемые. Имеется в виду суммарная длина отрезков
$\sum (y_i - x_i) < \delta$.
\end{warn}

\subsubsection*{Канторова лестница}

\inputfig{figures/fig-cantor}

Канторова лестница на самом канторовом множестве не задана явно, но продолжается
на весь отрезок по непрерывности. Она непрерывна, однако суммарная длина отрезков
$m$-го уровня равна $(2/3)^m \to 0$, тогда как суммарное приращение функции на них
остаётся равным $1$. Значит, \textbf{канторова лестница не абсолютно непрерывна}, хотя
«если абсолютно непрерывна, то непрерывна».

\begin{idea}
Это и показывает, что абсолютная непрерывность строго сильнее равномерной: именно она,
а не непрерывность, равносильна существованию суммарной производной.
\end{idea}

\begin{proposition}[эквивалентность определений]
Если $u$ абсолютно непрерывна в смысле определения через интеграл, то она абсолютно
непрерывна и в классическом смысле, и наоборот.
\end{proposition}

\begin{proof}
\plusplus
\end{proof}
